\chapter{Paires Supervisees et Fine-tuning}

\section{Construction des paires}
Le dataset de similarite final est construit sur des paires
\texttt{en <-> fr\_en} de la meme entite.
\begin{itemize}[leftmargin=2em]
  \item Positifs: meme entite, langues \texttt{en} et \texttt{fr\_en}.
  \item Negatifs: entites differentes (hard negatives + random negatives).
\end{itemize}

\begin{table}[H]
  \centering
  \caption{Split du dataset de paires \texttt{en <-> fr\_en} (hard final)}
  \begin{tabular}{lrrrr}
    \toprule
    Split & Lignes & Positifs & Negatifs & Entites \\
    \midrule
    Train & 1592 & 398 & 1194 & 199 \\
    Validation & 192 & 48 & 144 & 24 \\
    Test & 208 & 52 & 156 & 26 \\
    \bottomrule
  \end{tabular}
\end{table}

\section{Modeles compares}
\begin{itemize}[leftmargin=2em]
  \item Baseline non fine-tunee: \texttt{sentence-transformers/all-MiniLM-L6-v2}
  \item Fine-tune 1: \texttt{sidbrahim/narrativesAnalogues-allMiniLM}
  \item Fine-tune 2: \texttt{sidbrahim/narrativesAnalogues-MPNet}
\end{itemize}

\section{Objectif du fine-tuning}
L'objectif est d'ameliorer la separation entre paires analogues et non analogues,
en particulier sur les cas difficiles. La comparaison baseline vs modeles
fine-tunes permet de mesurer le gain reel de l'apprentissage supervise.
